\begin{table}[htbp]
\centering
\caption{Design-matrix multicollinearity diagnostics: pooled vs. modality-interaction GLM}
\label{tab:vif-matrix-summary}
\begin{tabular}{lcccccc}
\toprule
Model & Predictors & Full rank & Condition number & Max VIF & Terms VIF\textgreater5 & Terms VIF\textgreater10 \\
\midrule
Pooled & 5 & 12/12 rats & 3.86--4.13 & 4.17--4.71 & 0 & 0 \\
Modality interaction & 17 & 12/12 rats & 10.47--12.36 & 9.84--12.47 & 6--7 & 0--2 \\
\bottomrule
\end{tabular}
\vspace{4pt}
\par\noindent
\begin{minipage}{\linewidth}\small Range across the 12 rats. Condition number computed on the scaled design matrix (each column at unit variance), i.e. the ratio of largest to smallest singular value; a value below approximately 30 is conventionally not considered indicative of a numerically problematic design. VIF thresholds of 5 and 10 follow common convention. No infinite VIFs (i.e. no exact collinearity) were found in either model for any rat.\end{minipage}
\end{table}
